% GENERATED FILE. Edit canonical lessons, then run npm run generate:books.
% canonical-source-hash: fnv1a64:f3ddb5aded951bf7
% canonical-lessons: GE-C342-wirkung, GE-C342-wirken, GE-C342-verursachen, GE-C342-ausloesen, GE-C342-ausloeser

\chapter{Ursache und Wirkung}
\label{ch:ge-ursache-und-wirkung}
\hlchaptermodality{342}

\begin{chapteropening}
\textbf{By the end of this chapter:} \emph{I can talk about causes and effects: say what works, what causes a problem and what sets it off.}

Complete the last lesson of chapter 342: I can talk about causes and effects: say what works, what causes a problem and what sets it off.
\end{chapteropening}

\section[die Wirkung]{die Wirkung — the effect}

\label{lesson:GE-C342-wirkung}



\begin{quote}
Before the new one: say the German for nevertheless, then the German for in fact.
\end{quote}

\subsection*{You'll want to know: die Wirkung}
\textbf{die Wirkung} — \textquotedblleft{}the effect\textquotedblright{}.

\begin{quote}
Der Tee hat eine gute Wirkung: Ich schlafe besser.
\end{quote}

\begin{quote}
\emph{The tea has a good effect: I sleep better.}
\end{quote}

Plural: \emph{die Wirkungen}. \emph{Ursache und Wirkung} — cause and effect.

\begin{cousinweb}
\emph{Wirkung} is the noun of \emph{wirken}, inherited from Old High German \emph{wirken} and Proto-Germanic \emph{wirkijaną} \textquotedblleft{}to work, to make\textquotedblright{} (EtymDB), of the family of \emph{Werk} \textquotedblleft{}work\textquotedblright{}.
\end{cousinweb}

\begin{grammarlens}[title={What you've built}]
One more word for telling, explaining and describing at B1.
\end{grammarlens}

\subsection*{Your turn}
\begin{itemize}
\raggedright
  \item \emph{Say it:} \emph{die Wirkung}
  \item \emph{Say it:} \emph{die Wirkung}, once more
  \item \emph{Say it:} a sentence of your own with \emph{die Wirkung}
  \item \emph{Recall:} say the German for nevertheless, then the German for in fact, then say \emph{die Wirkung} again
\end{itemize}

\begin{tcolorbox}[breakable,colback=teal!4,colframe=teal!35!black,title={Before you move on}]
What does die Wirkung mean? (\textquotedblleft{}The effect\textquotedblright{}.) Say it once more.
\end{tcolorbox}
\section[wirken]{wirken — to work, to take effect; to seem}

\label{lesson:GE-C342-wirken}



\begin{quote}
Before the new one: say the German for in fact, then the German for the effect.
\end{quote}

\subsection*{You'll want to know: wirken}
\textbf{wirken} — \textquotedblleft{}to work, to take effect; to seem\textquotedblright{}.

\begin{quote}
Die Tabletten wirken schon nach einer halben Stunde.
\end{quote}

\begin{quote}
\emph{The tablets take effect after only half an hour.}
\end{quote}

It also means to seem: \emph{Er wirkt müde} — he seems tired.

\begin{cousinweb}
\emph{wirken} is inherited from Old High German \emph{wirken} and Proto-Germanic \emph{wirkijaną} \textquotedblleft{}to work, to make\textquotedblright{} (EtymDB), of the same family as \emph{Werk} \textquotedblleft{}work\textquotedblright{} and Greek \emph{érgon}.
\end{cousinweb}

\begin{grammarlens}[title={What you've built}]
One more word for telling, explaining and describing at B1.
\end{grammarlens}

\subsection*{Your turn}
\begin{itemize}
\raggedright
  \item \emph{Say it:} \emph{wirken}
  \item \emph{Say it:} \emph{wirken}, once more
  \item \emph{Say it:} a sentence of your own with \emph{wirken}
  \item \emph{Recall:} say the German for in fact, then the German for the effect, then say \emph{wirken} again
\end{itemize}

\begin{tcolorbox}[breakable,colback=teal!4,colframe=teal!35!black,title={Before you move on}]
What does wirken mean? (\textquotedblleft{}To work, to take effect; to seem\textquotedblright{}.) Say it once more.
\end{tcolorbox}
\section[verursachen]{verursachen — to cause}

\label{lesson:GE-C342-verursachen}



\begin{quote}
Before the new one: say the German for the effect, then the German for to work.
\end{quote}

\subsection*{You'll want to know: verursachen}
\textbf{verursachen} — \textquotedblleft{}to cause\textquotedblright{}.

\begin{quote}
Der starke Regen hat viele Probleme verursacht.
\end{quote}

\begin{quote}
\emph{The heavy rain caused a lot of problems.}
\end{quote}

Inseparable: \emph{es verursacht, es hat verursacht}.

\begin{cousinweb}
\emph{verursachen} is built with \emph{ver-} on \emph{Ursache} \textquotedblleft{}cause\textquotedblright{} (EtymDB), which joins the old prefix \emph{ur-} \textquotedblleft{}first, original\textquotedblright{} to \emph{Sache} \textquotedblleft{}thing, matter\textquotedblright{}, from Old High German \emph{sahha}.
\end{cousinweb}

\begin{grammarlens}[title={What you've built}]
One more word for telling, explaining and describing at B1.
\end{grammarlens}

\subsection*{Your turn}
\begin{itemize}
\raggedright
  \item \emph{Say it:} \emph{verursachen}
  \item \emph{Say it:} \emph{verursachen}, once more
  \item \emph{Say it:} a sentence of your own with \emph{verursachen}
  \item \emph{Recall:} say the German for the effect, then the German for to work, then say \emph{verursachen} again
\end{itemize}

\begin{tcolorbox}[breakable,colback=teal!4,colframe=teal!35!black,title={Before you move on}]
What does verursachen mean? (\textquotedblleft{}To cause\textquotedblright{}.) Say it once more.
\end{tcolorbox}
\section[auslösen]{auslösen — to trigger, to set off}

\label{lesson:GE-C342-ausloesen}



\begin{quote}
Before the new one: say the German for to work, then the German for to cause.
\end{quote}

\subsection*{You'll want to know: auslösen}
\textbf{auslösen} — \textquotedblleft{}to trigger, to set off\textquotedblright{}.

\begin{quote}
Seine Bemerkung hat einen großen Streit ausgelöst.
\end{quote}

\begin{quote}
\emph{His remark set off a big argument.}
\end{quote}

Separable: \emph{es löst aus}, \emph{es hat ausgelöst}. Also of machines: \emph{den Alarm auslösen} — to set off the alarm.

\begin{cousinweb}
\emph{auslösen} is built from \emph{aus} \textquotedblleft{}out\textquotedblright{} and \emph{lösen} \textquotedblleft{}to loosen, to release\textquotedblright{}, inherited from Old High German \emph{lōsen} and Proto-Germanic \emph{lausijaną}, from \emph{lausaz} \textquotedblleft{}loose\textquotedblright{} (EtymDB): you release something that then runs by itself.
\end{cousinweb}

\begin{grammarlens}[title={What you've built}]
One more word for telling, explaining and describing at B1.
\end{grammarlens}

\subsection*{Your turn}
\begin{itemize}
\raggedright
  \item \emph{Say it:} \emph{auslösen}
  \item \emph{Say it:} \emph{auslösen}, once more
  \item \emph{Say it:} a sentence of your own with \emph{auslösen}
  \item \emph{Recall:} say the German for to work, then the German for to cause, then say \emph{auslösen} again
\end{itemize}

\begin{tcolorbox}[breakable,colback=teal!4,colframe=teal!35!black,title={Before you move on}]
What does auslösen mean? (\textquotedblleft{}To trigger, to set off\textquotedblright{}.) Say it once more.
\end{tcolorbox}
\section[der Auslöser]{der Auslöser — the trigger}

\label{lesson:GE-C342-ausloeser}



\begin{quote}
Before the new one: say the German for to cause, then the German for to trigger.
\end{quote}

\subsection*{You'll want to know: der Auslöser}
\textbf{der Auslöser} — \textquotedblleft{}the trigger\textquotedblright{}.

\begin{quote}
Der Auslöser für den Streit war eine kleine Nachricht.
\end{quote}

\begin{quote}
\emph{The trigger for the argument was a small message.}
\end{quote}

Plural: \emph{die Auslöser}. On a camera, \emph{der Auslöser} is the shutter button.

\begin{cousinweb}
\emph{Auslöser} is formed from \emph{auslösen} \textquotedblleft{}to set off\textquotedblright{} with the ending \emph{-er}; the verb is built on \emph{lösen} \textquotedblleft{}to release\textquotedblright{}, from Proto-Germanic \emph{lausijaną} (EtymDB).
\end{cousinweb}

\begin{grammarlens}[title={What you've built}]
One more word for telling, explaining and describing at B1.
\end{grammarlens}

\subsection*{Your turn}
\begin{itemize}
\raggedright
  \item \emph{Say it:} \emph{der Auslöser}
  \item \emph{Say it:} \emph{der Auslöser}, once more
  \item \emph{Say it:} a sentence of your own with \emph{der Auslöser}
  \item \emph{Recall:} say the German for to cause, then the German for to trigger, then say \emph{der Auslöser} again
\end{itemize}

\begin{tcolorbox}[breakable,colback=teal!4,colframe=teal!35!black,title={Before you move on}]
What does der Auslöser mean? (\textquotedblleft{}The trigger\textquotedblright{}.) Say it once more.
\end{tcolorbox}
